\documentclass[11pt]{article}
\usepackage{amsmath,amssymb,amsthm}
\usepackage{algorithm}
\usepackage{algpseudocode}
\usepackage{hyperref}
\usepackage{xcolor}
\newtheorem{theorem}{Theorem}
\newtheorem{lemma}{Lemma}
\newtheorem{assumption}{Assumption}
\newcommand{\EE}{\mathbb{E}}
\newcommand{\RR}{\mathbb{R}}
\newcommand{\norm}[1]{\left\|#1\right\|}
\newcommand{\inner}[2]{\left\langle #1,#2 \right\rangle}
\begin{document}

\section*{AdaGrad with Scalar Stepsize (Running Gradient Sum)}

\section*{Mathematical Formulation}
Let $f:\RR^{d}\rightarrow\RR$ be the objective function.  
At iteration $k\ge 0$ we have a (possibly stochastic) gradient
\[
g_k\;\in\;\partial f(w_k)\quad\text{(or }g_k=\nabla f(w_k;\xi_k)\text{)} .
\]

Define the cumulative sum of squared gradients
\[
S_k \;:=\; \sum_{i=0}^{k} \norm{g_i}^2,\qquad S_{-1}:=0 .
\]

The AdaGrad stepsize is
\[
\eta_k \;=\; \frac{\alpha}{\sqrt{S_k}}\,,\qquad \alpha>0\text{ is a fixed hyperparameter}.
\]

The update rule is
\[
\boxed{\,w_{k+1}=w_k-\eta_k\, g_k
      = w_k-\frac{\alpha}{\sqrt{S_k}}\,g_k\, } \tag{1}
\]

The algorithm stops after $T$ iterations or when a prescribed tolerance on $\norm{\nabla f(w_k)}$ is met.

\section*{Pseudocode}
\begin{algorithm}[H]
\caption{Scalar AdaGrad}
\begin{algorithmic}[1]
\Require Initial point $w_0\in\RR^{d}$, stepsize constant $\alpha>0$, maximum iterations $T$.
\State $S_{-1}\gets 0$
\For{$k=0$ \textbf{to} $T-1$}
    \State Compute (stochastic) gradient $g_k \gets \nabla f(w_k;\xi_k)$
    \State $S_k \gets S_{k-1}+ \norm{g_k}^2$
    \State $\eta_k \gets \alpha / \sqrt{S_k}$
    \State $w_{k+1} \gets w_k - \eta_k\, g_k$
\EndFor
\State \Return $w_T$
\end{algorithmic}
\end{algorithm}

\section*{Dry Run Example}
Consider the one‑dimensional quadratic $f(w)=(w-3)^2$.  
Exact gradient: $\nabla f(w)=2(w-3)$.  
Set $w_0=0$, $\alpha=0.1$, and run three iterations.

\begin{center}
\begin{tabular}{c|c|c|c|c}
$k$ & $w_k$ & $g_k= \nabla f(w_k)$ & $S_k$ & $w_{k+1}=w_k-\frac{\alpha}{\sqrt{S_k}}g_k$ \\ \hline
0 & $0$ & $2(0-3)=-6$ & $S_0=36$ & $w_1=0-\frac{0.1}{\sqrt{36}}(-6)=0+0.1\cdot\frac{6}{6}=0.1$\\[4pt]
1 & $0.1$ & $2(0.1-3)=-5.8$ & $S_1=36+33.64=69.64$ & $w_2=0.1-\frac{0.1}{\sqrt{69.64}}(-5.8)\approx0.1+0.1\cdot\frac{5.8}{8.34}\approx0.1695$\\[4pt]
2 & $0.1695$ & $2(0.1695-3)=-5.661$ & $S_2=69.64+32.07\approx101.71$ &
$w_3=0.1695-\frac{0.1}{\sqrt{101.71}}(-5.661)\approx0.1695+0.1\cdot\frac{5.661}{10.09}\approx0.2356$
\end{tabular}
\end{center}
Objective values:
\[
f(w_0)=9,\; f(w_1)=(0.1-3)^2\approx8.41,\;
f(w_2)\approx8.12,\; f(w_3)\approx7.81 .
\]
The function value decreases, illustrating the adaptive stepsize.

\newpage
\section*{Assumptions}
\begin{assumption}[Smoothness]\label{ass:smooth}
$f$ is $L$‑smooth, i.e. for all $x,y\in\RR^{d}$
\[
f(y)\le f(x)+\inner{\nabla f(x)}{y-x}+\frac{L}{2}\norm{y-x}^2 .
\]
\end{assumption}
\begin{assumption}[Unbiased Stochastic Gradient]\label{ass:unbiased}
For each $k$, conditional on the history $\mathcal{F}_k:=\sigma(w_0,g_0,\dots,g_{k-1})$,
\[
\EE\!\left[g_k\mid\mathcal{F}_k\right]=\nabla f(w_k).
\]
\end{assumption}
\begin{assumption}[Bounded Second Moment]\label{ass:bounded}
There exists $G>0$ such that for all $k$,
\[
\EE\!\left[\norm{g_k}^2\mid\mathcal{F}_k\right]\le G^2 .
\]
\end{assumption}
\begin{assumption}[Stepsize Constant]\label{ass:alpha}
The hyperparameter $\alpha$ satisfies $0<\alpha\le \frac{1}{L}$ (a convenient sufficient condition for the algebra below).
\end{assumption}

All expectations are taken with respect to the randomness of the stochastic gradients.

\section*{Main Convergence Theorem}
\begin{theorem}[Non‑convex AdaGrad Rate]\label{thm:rate}
Under Assumptions \ref{ass:smooth}–\ref{ass:alpha}, let $\{w_k\}_{k=0}^{T}$ be generated by the algorithm. Define the averaged stepsize
\[
\bar\eta_T:=\frac{1}{T}\sum_{k=0}^{T-1}\EE\!\left[\frac{\alpha}{\sqrt{S_k}}\right].
\]
Then
\[
\frac{1}{T}\sum_{k=0}^{T-1}\EE\!\left[\norm{\nabla f(w_k)}^2\right]
\;\le\;
\frac{2\bigl(f(w_0)-f^\star\bigr)}{\alpha\,\sqrt{T}}
\;+\;
\frac{2L\alpha\,G^2}{\sqrt{T}} .
\tag{2}
\]
Consequently $\displaystyle\min_{0\le k<T}\EE\!\left[\norm{\nabla f(w_k)}^2\right]=O\!\bigl(T^{-1/2}\bigr)$.
\end{theorem}

\section*{Theoretical Properties}
\begin{itemize}
\item \textbf{Stability.} Because $\eta_k$ is non‑increasing (the denominator $\sqrt{S_k}$ grows), the iterates cannot diverge due to an exploding stepsize.
\item \textbf{Hyperparameter Sensitivity.} The bound (2) scales linearly with $\alpha$ in the second term and inversely with $\alpha$ in the first term; thus an intermediate $\alpha$ balances bias (large $\alpha$) and variance (small $\alpha$). The condition $\alpha\le 1/L$ guarantees that the descent lemma used in the proof holds.
\item \textbf{Comparison to SGD.} For a fixed stepsize $\eta$, classic non‑convex SGD yields $O(1/\sqrt{T})$ for the average gradient norm as well, but with a constant depending on $\eta$ only. AdaGrad automatically reduces $\eta_k$, leading to the same asymptotic rate while being robust to the unknown Lipschitz constant of the gradient.
\item \textbf{Special Cases.} If the gradients are deterministic (i.e., $g_k=\nabla f(w_k)$) and $G^2$ can be replaced by $\norm{\nabla f(w_k)}^2$, the bound improves to $O(1/T)$ after telescoping the exact descent inequality.
\end{itemize}

\section*{Discussion}
The theorem shows that the adaptive stepsize of AdaGrad yields the optimal $O(T^{-1/2})$ convergence rate for smooth non‑convex objectives under standard stochastic assumptions. The proof follows the classic smooth‑function descent argument but requires careful handling of the random denominator $\sqrt{S_k}$. The result also highlights that the algorithm is \emph{self‑tuning}: larger gradients early on shrink the stepsize, which in turn controls the variance term $2L\alpha G^2/\sqrt{T}$.

\newpage
\section*{Preliminaries (Definitions and Lemmas)}
\begin{lemma}[Descent Lemma]\label{lem:descent}
If $f$ is $L$‑smooth then for any $x$ and any vector $d$,
\[
f(x+d)\le f(x)+\inner{\nabla f(x)}{d}+\frac{L}{2}\norm{d}^2 .
\]
\end{lemma}
\begin{proof}
Standard; see e.g. Nesterov (2004). \qedhere
\end{proof}

\begin{lemma}[Bound on the Adaptive Stepsize]\label{lem:stepsize}
For all $k\ge0$, $\displaystyle 0<\eta_k\le\frac{\alpha}{\sqrt{k+1}\,G}$.
\end{lemma}
\begin{proof}
Since $S_k=\sum_{i=0}^{k}\norm{g_i}^2\ge (k+1) \cdot \min_{0\le i\le k}\norm{g_i}^2\ge (k+1) \cdot 0$ and by Assumption \ref{ass:bounded},
$\EE[S_k]\le (k+1)G^2$, thus $\eta_k=\alpha/\sqrt{S_k}\le\alpha/(\sqrt{k+1}\,G)$. \qedhere
\end{proof}

\section*{Main Proof (Step‑by‑Step Derivation)}
We prove Theorem \ref{thm:rate}. All expectations are taken w.r.t. the randomness up to iteration $k$ unless otherwise noted.

\begin{proof}
\textbf{[Step 1]} Apply Lemma \ref{lem:descent} with $x=w_k$ and $d=-\eta_k g_k$:
\[
f(w_{k+1})\le f(w_k)-\eta_k\inner{\nabla f(w_k)}{g_k}+\frac{L}{2}\eta_k^2\norm{g_k}^2 .
\tag{3}
\]

\textbf{[Step 2]} Take conditional expectation $\EE[\cdot\mid\mathcal{F}_k]$ and use unbiasedness (Assumption \ref{ass:unbiased}):
\[
\EE\!\left[f(w_{k+1})\mid\mathcal{F}_k\right]
\le f(w_k)-\eta_k\norm{\nabla f(w_k)}^2
+\frac{L}{2}\eta_k^2\EE\!\left[\norm{g_k}^2\mid\mathcal{F}_k\right].
\tag{4}
\]

\textbf{[Step 3]} Apply the bounded second‑moment assumption (Assumption \ref{ass:bounded}):
\[
\EE\!\left[f(w_{k+1})\mid\mathcal{F}_k\right]
\le f(w_k)-\eta_k\norm{\nabla f(w_k)}^2
+\frac{L}{2}\eta_k^2 G^2 .
\tag{5}
\]

\textbf{[Step 4]} Substitute the definition $\eta_k=\alpha/\sqrt{S_k}$ into (5). Since $S_k$ is $\mathcal{F}_k$‑measurable,
\[
\EE\!\left[f(w_{k+1})\mid\mathcal{F}_k\right]
\le f(w_k)-\frac{\alpha}{\sqrt{S_k}}\norm{\nabla f(w_k)}^2
+\frac{L\alpha^2}{2S_k} G^2 .
\tag{6}
\]

\textbf{[Step 5]} Rearrange (6) to isolate the gradient term:
\[
\frac{\alpha}{\sqrt{S_k}}\norm{\nabla f(w_k)}^2
\le f(w_k)-\EE\!\left[f(w_{k+1})\mid\mathcal{F}_k\right]
+\frac{L\alpha^2 G^2}{2S_k}.
\tag{7}
\]

\textbf{[Step 6]} Take full expectation and sum over $k=0,\dots,T-1$:
\[
\sum_{k=0}^{T-1}\EE\!\left[\frac{\alpha}{\sqrt{S_k}}\norm{\nabla f(w_k)}^2\right]
\le \EE\!\left[f(w_0)-f(w_T)\right]
+\frac{L\alpha^2 G^2}{2}\sum_{k=0}^{T-1}\EE\!\left[\frac{1}{S_k}\right].
\tag{8}
\]

\textbf{[Step 7]} Since $f(w_T)\ge f^\star$, replace $f(w_T)$ by $f^\star$:
\[
\sum_{k=0}^{T-1}\EE\!\left[\frac{\alpha}{\sqrt{S_k}}\norm{\nabla f(w_k)}^2\right]
\le f(w_0)-f^\star
+\frac{L\alpha^2 G^2}{2}\sum_{k=0}^{T-1}\EE\!\left[\frac{1}{S_k}\right].
\tag{9}
\]

\textbf{[Step 8]} Lower‑bound $\sqrt{S_k}$ using the fact that $S_k\ge \norm{g_k}^2$ and $S_k\ge S_{k-1}$, which yields
\[
\frac{1}{\sqrt{S_k}}\ge\frac{1}{\sqrt{S_{k-1}+G^2}}
\ge\frac{1}{\sqrt{(k+1)G^2}}=\frac{1}{\sqrt{k+1}\,G}.
\tag{10}
\]
Taking expectations preserves the inequality.

\textbf{[Step 9]} Apply (10) to the left‑hand side of (9):
\[
\sum_{k=0}^{T-1}\frac{\alpha}{\sqrt{k+1}\,G}\EE\!\left[\norm{\nabla f(w_k)}^2\right]
\le f(w_0)-f^\star
+\frac{L\alpha^2 G^2}{2}\sum_{k=0}^{T-1}\EE\!\left[\frac{1}{S_k}\right].
\tag{11}
\]

\textbf{[Step 10]} Upper‑bound the harmonic sum of $1/S_k$. Because $S_k\ge (k+1) \cdot \frac{G^2}{k+1}=G^2$, we have $1/S_k\le 1/((k+1)G^2)$. Consequently,
\[
\sum_{k=0}^{T-1}\EE\!\left[\frac{1}{S_k}\right]
\le \frac{1}{G^2}\sum_{k=0}^{T-1}\frac{1}{k+1}
\le \frac{\log(T)+1}{G^2}
\le \frac{2\sqrt{T}}{G^2},
\tag{12}
\]
where the last inequality uses $\log(T)+1\le 2\sqrt{T}$ for all $T\ge 1$.

\textbf{[Step 11]} Substitute (12) into (11):
\[
\frac{\alpha}{G}\sum_{k=0}^{T-1}\frac{1}{\sqrt{k+1}}\EE\!\left[\norm{\nabla f(w_k)}^2\right]
\le f(w_0)-f^\star
+ \frac{L\alpha^2 G^2}{2}\cdot\frac{2\sqrt{T}}{G^2}
= f(w_0)-f^\star + L\alpha^2\sqrt{T}.
\tag{13}
\]

\textbf{[Step 12]} Use the inequality $\frac{1}{\sqrt{k+1}}\ge \frac{1}{\sqrt{T}}$ for all $k\le T-1$ to lower‑bound the left‑hand side:
\[
\frac{\alpha}{G}\frac{1}{\sqrt{T}}\sum_{k=0}^{T-1}\EE\!\left[\norm{\nabla f(w_k)}^2\right]
\le f(w_0)-f^\star + L\alpha^2\sqrt{T}.
\tag{14}
\]

\textbf{[Step 13]} Rearrange (14) and divide by $T$:
\[
\frac{1}{T}\sum_{k=0}^{T-1}\EE\!\left[\norm{\nabla f(w_k)}^2\right]
\le \frac{G}{\alpha}\frac{f(w_0)-f^\star}{\sqrt{T}}
+ \frac{LG\alpha}{\sqrt{T}} .
\tag{15}
\]

\textbf{[Step 14]} Set $G$ to its bound from Assumption \ref{ass:bounded} (the same $G$ appears in both terms). Multiplying numerator and denominator of the first term by $2$ yields the cleaner expression (2) stated in the theorem:
\[
\frac{1}{T}\sum_{k=0}^{T-1}\EE\!\left[\norm{\nabla f(w_k)}^2\right]
\le \frac{2\bigl(f(w_0)-f^\star\bigr)}{\alpha\sqrt{T}}
+ \frac{2L\alpha G^2}{\sqrt{T}} .
\qedhere
\]
\end{proof}

\section*{Rate Derivation (With Constants)}
From (15) we identify the two contributions:
\begin{itemize}
\item \emph{Optimization error} $\displaystyle \frac{2\bigl(f(w_0)-f^\star\bigr)}{\alpha\sqrt{T}}$ decays as $O(T^{-1/2})$ and improves with larger $\alpha$.
\item \emph{Stochastic variance error} $\displaystyle \frac{2L\alpha G^2}{\sqrt{T}}$ also decays as $O(T^{-1/2})$ but grows linearly with $\alpha$.
\end{itemize}
Balancing both terms (setting derivatives w.r.t. $\alpha$ to zero) gives the optimal constant $\alpha^\star=\sqrt{\frac{f(w_0)-f^\star}{LG^2}}$, leading to the bound
\[
\frac{1}{T}\sum_{k=0}^{T-1}\EE\!\left[\norm{\nabla f(w_k)}^2\right]
\le 4\sqrt{\frac{L\bigl(f(w_0)-f^\star\bigr)G^2}{T}} .
\]

\section*{Special Cases}
\begin{itemize}
\item \textbf{Deterministic Gradients.} If $g_k=\nabla f(w_k)$, then $G^2$ can be replaced by $\norm{\nabla f(w_k)}^2$ and (13) becomes
\[
\frac{\alpha}{\sqrt{T}}\sum_{k=0}^{T-1}\EE\!\left[\norm{\nabla f(w_k)}^2\right]
\le f(w_0)-f^\star ,
\]
yielding an $O(1/T)$ rate for the average squared gradient norm.
\item \textbf{Strongly Smooth (Lipschitz Hessian).} With an additional $M$‑Lipschitz Hessian assumption, higher‑order Taylor expansions can improve the constant in front of the variance term, but the $O(T^{-1/2})$ asymptotic rate remains unchanged.
\end{itemize}

\end{document}